\documentclass[9pt,a4paper]{article}
\usepackage[a4paper,margin=1.2cm,top=1.2cm,bottom=1.2cm]{geometry}
\usepackage{fontspec}
\usepackage[brazilian]{babel}

% --- TIPOGRAFIA (NEJM Style) ---
\setmainfont{Charter}[
  Scale=1.0,
  Ligatures=TeX,
  Contextuals=Alternate
]
\setsansfont{Helvetica Neue}[Scale=0.9]
\newfontfamily\titlefont{Helvetica Neue Condensed Bold}[Scale=1.1]
\newfontfamily\headerfont{Helvetica Neue Bold}[Scale=0.9]

% --- PACOTES ---
\usepackage{microtype} % Justificação e kerning superiores
\usepackage{xcolor}
\usepackage{booktabs}
\usepackage{multicol}
\usepackage{amsmath}
\usepackage{tikz}
\usetikzlibrary{calc, positioning, arrows.meta}

% --- CORES NEJM ---
\definecolor{nejmBlue}{RGB}{0, 75, 145}
\definecolor{nejmGray}{RGB}{80, 80, 80}
\definecolor{nejmLightGray}{RGB}{245, 245, 245}

\pagestyle{empty}

% --- AJUSTES DE ESPAÇAMENTO ---
\setlength{\columnsep}{0.6cm}
\setlength{\columnseprule}{0.4pt}
\renewcommand{\columnseprulecolor}{\color{nejmLightGray}}

% --- MACROS ---
\newcommand{\dropcap}[1]{{\color{nejmBlue}\fontsize{24}{28}\selectfont\bfseries #1}}

\newcommand{\quicktake}[1]{%
    \begin{center}
    \fcolorbox{nejmBlue}{nejmLightGray}{%
        \begin{minipage}{\dimexpr\linewidth-10pt}
            \vspace{2pt}
            {\headerfont\small\textcolor{nejmBlue}{QUICK TAKE}}\\[2pt]
            {\scriptsize #1}
            \vspace{2pt}
        \end{minipage}%
    }
    \end{center}
}

\newcommand{\sectionheader}[1]{%
    \vspace{0.3cm}
    \noindent{\headerfont\small\textcolor{nejmBlue}{\MakeUppercase{#1}}}%
    \par\vspace{1pt}\hrule height 0.8pt color nejmBlue
    \vspace{0.15cm}
}

\begin{document}

% === HEADER NEJM STYLE ===
\noindent
\begin{minipage}[b]{0.7\textwidth}
    {\fontsize{8}{10}\selectfont\headerfont\textcolor{nejmGray}{CLINICAL PRACTICE}} \\[4pt]
    {\titlefont\fontsize{22}{24}\selectfont\textcolor{black}{Pulse Oximetry in Critical Care}}
\end{minipage}\hfill
\begin{minipage}[b]{0.28\textwidth}
    \begin{flushright}
        {\fontsize{7}{9}\selectfont\headerfont\textcolor{nejmBlue}{ONE PAGER ICU}} \\
        {\fontsize{7}{9}\selectfont\textit{Tradução Técnica e Adaptação}} \\
        {\fontsize{7}{9}\selectfont\textcolor{nejmGray}{2025 Edition}}
    \end{flushright}
\end{minipage}

\vspace{0.2cm}
\hrule height 2pt color nejmBlue
\vspace{0.1cm}
\hrule height 0.5pt color nejmGray
\vspace{0.4cm}

\begin{multicols}{2}

\dropcap{P}ulse oximetry is a ubiquitous noninvasive technology used for continuous monitoring of arterial oxygen saturation ($SpO_2$). In the intensive care unit (ICU), it serves as a critical surrogate for arterial blood gas analysis ($SaO_2$), enabling rapid detection of hypoxemia. However, the apparent simplicity of the technology masks complex physical principles and significant clinical limitations that can lead to diagnostic errors.

\sectionheader{Physiological Principles}
The technology relies on \textbf{spectrophotometry} and \textbf{plethysmography}. It exploits the fact that oxygenated hemoglobin ($HbO_2$) and deoxygenated hemoglobin ($Hb$) have distinct absorption spectra. 

\textit{Beer-Lambert Law}: The amount of light absorbed is proportional to the concentration of the absorbing substance and the path length. Pulse oximeters use two specific wavelengths:
\begin{itemize}
    \setlength{\itemsep}{0pt}
    \item \textbf{Red (660\,nm)}: Highly absorbed by $Hb$.
    \item \textbf{Infrared (940\,nm)}: Highly absorbed by $HbO_2$.
\end{itemize}

\textit{Plethysmography}: To isolate arterial saturation, the device identifies the \textbf{pulsatile component} (AC) of the signal, which represents the varying volume of arterial blood, and filters out the constant absorption (DC) from tissue, venous blood, and bone.

\vspace{0.3cm}
\begin{center}
\begin{tikzpicture}[scale=0.8, font=\sffamily\tiny]
    \draw[->, nejmGray] (0,0) -- (5.5,0) node[right, nejmBlue] {$\lambda$ (nm)};
    \draw[->, nejmGray] (0,0) -- (0,3.5) node[above, nejmBlue] {Absorp.};
    
    % Curvas (Nature/NEJM style)
    \draw[thick, red!80!black] (0.5, 3.2) .. controls (1.5, 2.5) and (3, 1) .. (5, 0.5) node[right] {$HbO_2$};
    \draw[thick, blue!80!black] (0.5, 0.4) .. controls (1.5, 1.2) and (3, 2.8) .. (5, 2.8) node[right] {$Hb$};
    
    % LED bars
    \fill[red, opacity=0.1] (0.9, 0) rectangle (1.1, 3.5);
    \draw[red!50, dashed] (1.0, 0) -- (1.0, 3.5) node[above] {660};
    \fill[orange, opacity=0.1] (4.4, 0) rectangle (4.6, 3.5);
    \draw[orange!50, dashed] (4.5, 0) -- (4.5, 3.5) node[above] {940};
\end{tikzpicture}
\end{center}

\quicktake{Pulse oximetry measures $SpO_2$, not $PaO_2$. In hyperoxemia ($SpO_2 = 100\%$), the device cannot detect dangerously high oxygen partial pressures (hyperoxia), which may contribute to oxidative stress.}

\sectionheader{Clinical Accuracy and Bias}
The accuracy of $SpO_2$ is typically $\pm 2\%$ when saturations are above 80\%. Below 70\%, the correlation with $SaO_2$ deteriorates significantly as the device's empirical calibration curves are derived from healthy volunteers.

\textbf{Skin Pigmentation and Racial Bias}: 
Multiple studies have demonstrated that pulse oximetry may \textbf{overestimate} $SaO_2$ in patients with dark skin. This bias leads to "occult hypoxemia," where patients are triaged as stable despite real saturations being dangerously low. Black patients face a 3-fold higher risk of this diagnostic error compared to white patients.

\textbf{Low Perfusion States}: 
In states of shock, severe vasoconstriction, or hypothermia, the pulsatile signal may be too weak to detect. The \textbf{Perfusion Index (PI)} can quantify this; a PI below 0.4\% often indicates unreliable data.

\sectionheader{Abnormal Hemoglobins}
\begin{itemize}
    \setlength{\itemsep}{0pt}
    \item \textbf{Carboxyhemoglobin ($COHb$)}: CO has a similar absorption to $HbO_2$ at 660\,nm. Consequently, $SpO_2$ remains near 100\% even in lethal CO poisoning.
    \item \textbf{Methemoglobin ($MetHb$)}: Distorts absorption at both wavelengths, usually resulting in a "plateau" reading around 85\%.
\end{itemize}

\sectionheader{Waveform Analysis}
Verification of the plethysmographic waveform is essential. A high-quality wave should show clear systolic and diastolic phases.

\vspace{0.2cm}
\begin{center}
\begin{tikzpicture}[scale=0.9, font=\sffamily\tiny]
    \draw[->, nejmGray] (0,0) -- (4.5,0);
    \draw[thick, nejmBlue] (0.2, 0.1) .. controls (0.4, 0.1) and (0.5, 1.4) .. (0.7, 1.4) 
         .. controls (0.9, 1.4) and (1.0, 0.6) .. (1.1, 0.7) 
         .. controls (1.2, 0.8) and (1.6, 0.1) .. (2.0, 0.1)
         .. controls (2.2, 0.1) and (2.3, 1.4) .. (2.5, 1.4)
         .. controls (2.7, 1.4) and (2.8, 0.6) .. (2.9, 0.7)
         .. controls (3.0, 0.8) and (3.4, 0.1) .. (3.8, 0.1);
    \node[anchor=south, nejmBlue, font=\bfseries] at (0.7, 1.4) {Systole};
    \node[anchor=west, nejmGray] at (1.1, 0.8) {Dicrotic Notch};
\end{tikzpicture}
\end{center}

\sectionheader{Interferences and Technical Pitfalls}
External factors like nail polish (especially black, blue, and green) and bright ambient lighting can saturate the detector. For patients with suspected circulatory failure, a probe placed on the forehead or earlobe may provide more rapid and reliable data than finger probes.

\sectionheader{Final Considerations \& Clinical Pearls}
\begin{itemize}
    \setlength{\itemsep}{0pt}
    \item \textbf{Ear Loop vs. Finger}: The earloop probe is less susceptible to vasoconstriction-induced error.
    \item \textbf{Pulse rate vs. EKG}: Always ensure the plethysmographic pulse rate matches the EKG heart rate.
    \item \textbf{Methylene Blue}: Can cause a precipitate drop in $SpO_2$ to 65\% within 1--2 minutes of administration.
\end{itemize}

\vspace{0.2cm}
\begin{center}
\begin{tabular}{lp{4.2cm}}
    \toprule
    \textbf{Factor} & \textbf{Clinical Impact on $SpO_2$} \\
    \midrule
    $COHb$ & Falsely normal; masks hypoxia \\
    $MetHb$ & Fixed reading near $\sim$85\% \\
    Skin Tone & Potential overestimation ($>$ 4\%) \\
    Shock & Signal loss or instability \\
    Methylene Blue & Spurious sudden drop \\
    \bottomrule
\end{tabular}
\end{center}

\vspace{0.2cm}
\noindent\textit{Authoritative Note}: When $SpO_2$ values conflict with clinical observation, \textbf{arterial blood gas analysis with co-oximetry} remains the mandatory gold standard for definitive diagnosis.

\end{multicols}

\vspace{0.4cm}
\hrule height 0.5pt
\vspace{1pt}
\noindent {\fontsize{6}{8}\selectfont \textcolor{nejmGray}{Copyright \copyright\ 2025 ICU Reference Group. All rights reserved. Source material adapted from Nick Mark, MD, for academic dissemination.}}

\end{document}
